\begin{textsample}{POS Dim 5 – global\_south\_2024\_12 – Score 21.00 – global\_south\_2024\_12/119054668.txt}  \label{ex:f5_pos_041}
A Palestinian child carries a water container, as Palestinians wait to collect water following a Human Rights Watch report that says Israel ' s deprivation of water in Gaza is an \textbf{act} of \textbf{genocide} in Khan Younis in the southern Gaza Strip December 19. -- AFP Human Rights Watch ( HRW ) on Thursday said Israel was \textbf{committing}"\textbf{acts} of \textbf{genocide}"in the Gaza Strip by damaging water infrastructure and cutting off supplies to civilians, calling on the \textbf{international} community to impose targeted \textbf{sanctions}.
In a new report, which focused specifically on water, the New York-based rights group detailed what it said were deliberate efforts by Israeli authorities"of a systematic nature"to deprive Palestinians in Gaza of water, which had"likely caused thousands of deaths... and will likely continue to cause deaths"."Since October 2023, Israeli authorities have deliberately obstructed Palestinians ' access to the adequate amount of water required for survival in the Gaza Strip,"the report said.
Israel on Thursday said the report by Human Rights Watch accusing it of \textbf{committing}"\textbf{acts} of \textbf{genocide}"in the Gaza Strip by restricting access to water for civilians was"full of lies".
The HRW report detailed what the group said was the intentional damaging of water and sanitation infrastructure, including solar panels powering treatment plants, a reservoir and a spare parts warehouse, as well as the blocking of fuel for generators.
The report concluded that in doing so,"Israeli authorities \textbf{intentionally} inflicted on the Palestinian population in Gaza ' conditions of life calculated to bring about its physical destruction in whole or in part. '"This, it said, amounted to the war \textbf{crime} of"extermination"and to"\textbf{acts} of \textbf{genocide}".
Speaking to Dawn.com on the report, HRW Senior Counsel Asia Division Saroop Ijaz said,"This is extermination, a \textbf{crime} against \textbf{humanity}, and an \textbf{act} of \textbf{genocide}."He said that Israeli authorities didn't just cut off water for Palestinians in Gaza.
They also destroyed water infrastructure, targeted water workers, and blocked water aid."This very deliberate \textbf{pattern} of conduct is ' calculated ( for ) physical destruction '.
That makes it an \textbf{act} of \textbf{genocide}."He also added,"governments arming Israel and undermining accountability efforts risk complicity in Israel ' s atrocities in Gaza -- Countries have a \textbf{legal} and moral responsibility to prevent \textbf{genocide} against Palestinians by arms \textbf{embargo}, targeted \textbf{sanctions} and support for justice."Human Rights Watch is the second major rights group in a month to use the word \textbf{genocide} to describe the actions of Israel in Gaza after \textbf{Amnesty} \textbf{International} issued a report that concluded Israel was \textbf{committing} \textbf{genocide}.
A report from Doctors Without Borders ( MSF ) released today has also found"clear signs of \textbf{ethnic} \textbf{cleansing}"in Gaza by Israeli attacks However, HRW stopped short of saying Israel was \textbf{committing} outright"\textbf{genocide}".
Under \textbf{international} law, proving \textbf{genocide} requires evidence of specific \textbf{intent}, which experts say is very difficult.
HRW said that,"The \textbf{pattern} of conduct set out in this report together with statements suggesting some Israeli officials wished to destroy the Palestinians in Gaza may indicate such \textbf{intent}."Speaking at a briefing on the report, Lama Faqih, director of HRW ' s Middle East and North Africa division, said that in the absence of"a clear articulated plan"to \textbf{commit} \textbf{genocide}, the \textbf{International} \textbf{Court} of Justice ( ICJ ) might find that the evidence meets the"very strict threshold"of \textbf{reasonable} inference of \textbf{genocidal} \textbf{intent}.
HRW pointed to a statement by then-defence minister Yoav Gallant in October 2023, when he declared a"complete siege"and said :"No electricity, no food, no water, no gas -- it ' s all closed."Israel is facing a \textbf{case} brought by South Africa at the ICJ last December, arguing that Israel ' s invasion of Gaza breached the 1948 United Nations \textbf{Genocide} \textbf{Convention}, an accusation Israel has strongly denied.
Israel denies intentional destruction of the population in Gaza, saying that it facilitates aid to the besieged territory. ' Malnourished and dehydrated ' The HRW report, drawn up over nearly a year, is based on interviews with dozens of Palestinians in Gaza, staff at water and sanitation facilities, medics and aid workers, as well as satellite imagery, photographs, videos and data \textbf{analysis}.
It said Israeli authorities did not reply to requests for information.
The lack of water left Palestinians in Gaza vulnerable to water-borne diseases and complications, such as infected wounds and the inability to heal due to dehydration, HRW said.

% matched lemmas: act, amnesty, analysis, case, cleansing, commit, convention, court, crime, embargo, ethnic, genocidal, genocide, humanity, intent, intentionally, international, legal, pattern, reasonable, sanction
\end{textsample}
